% !PRINTED 267
\weberpage{267}
%<<CONT
set, it can be represented as follows:
\weq{9}{\Phi_p(v,u)=\Bigl(\frac{uv}{\sqrt2}\Bigr)^{p+1}\Phi_p\Bigl(\frac{\varepsilon\sqrt2}{v},\frac{\sqrt2}{u}\Bigr).}
From this one infers the relations
\weq{10}{\varepsilon^h2^{\frac{h+k-p-1}{2}}a_{p+1-h,\,p+1-k}=a_{h,k}, \qquad a_{h,k}=a_{k,h}.}
If, therefore, we collect together in (7) the terms with equal
coefficients $a_{h,k}$, we may restrict ourselves to the assumption
that $h\geqq k$, $h+k\geqq p+1$; and, apart from the two terms
$v^{p+1}, u^{p+1}$, $\Phi_p$ presents itself as an aggregate of
terms of the two following forms (if it is further taken into
account that, on account of (8), $\varepsilon^h=\varepsilon^k$):
\weq{11}{\begin{aligned}
&v^hu^k+v^ku^h+\varepsilon^h2^{\frac{h+k-p-1}{2}}\bigl(u^{p+1-h}v^{p+1-k}+u^{p+1-k}v^{p+1-h}\bigr)=\\
&\underset{h>k}{(uv)^{\frac{p+1}{2}}}\left[\Bigl(\frac{v}{u}\Bigr)^{\frac{h-k}{2}}+\Bigl(\frac{u}{v}\Bigr)^{\frac{h-k}{2}}\right]\left[(uv)^{\frac{h+k-p-1}{2}}+\frac{\varepsilon^h2^{\frac{h+k-p-1}{2}}}{(uv)^{\frac{h+k-p-1}{2}}}\right].
\end{aligned}}
\weq{12}{\begin{aligned}
&v^hu^h+\varepsilon^h2^{h-\frac{p+1}{2}}(uv)^{p+1-h}\\
&=(uv)^{\frac{p+1}{2}}\left[(uv)^{h-\frac{p+1}{2}}+\frac{\varepsilon^h2^{h-\frac{p+1}{2}}}{(uv)^{h-\frac{p+1}{2}}}\right].
\end{aligned}}
The coefficients of these terms in $\Phi_p$ are integers.
Now from (8) there results:
\[
\begin{aligned}
(h-k)&\equiv (h-1)(p+1)\\
h+k-p-1&\equiv -h(p-1)
\end{aligned}
\ (\mathrm{mod}\ 24),
\]
and therefore $h-k$ is divisible by $2r$, and $h+k-p-1$ (in which
$h$ may also be $=k$) by $2s$ [\S\ 74, (2)].

If we therefore set
\[
\frac{h-k}{2}=r\alpha, \qquad \frac{h+k-p-1}{2}=s\beta,
\]
then $\alpha$ and $\beta$ are positive integers which lie between
the limits
\weq{13}{0\leqq\alpha<\frac{p+1}{2r}, \qquad 0\leqq\beta\leqq\frac{p+1}{2s},}
and moreover the congruence $-h(p-1)\equiv 2s\beta\ (\mathrm{mod}\ 24)$,
which follows from (8), shows that at least in the cases where
$\varepsilon=-1$, $h\equiv\beta\ (\mathrm{mod}\ 2)$
[\S\ 74, (2)], and hence always $\varepsilon^h=\varepsilon^\beta$.
